\clearpage
\subsection{An Example of a Topic that involves Varieties and Polynomials: Parametrization and Implicitization} 

\begin{exercise}[title = Example 1.3.1 (Linear Case)]
  Let's consider a line, given implicitly, that means by equations. 
  \begin{align*}
  \left\{\begin{aligned}
    x + y + z & = 1 \\
    x + 2y -z & = 3
  \end{aligned}\right. 
\end{align*}
Our field is $\mathbb{R}$ for this example. Describe this line parametically. Let's call the parameter $t$.  \\ 

\hspace*{0.5cm}\textcolor{airforceblue}{\textit{(Flashback to Kindergarten!!)}} \\ 

\paragraph{Solution:} $\rightarrow$ same setting as linear case in Example 1.2.2 (6.)
\begin{align*}
  \left\{
  \begin{aligned}
    x + y + z & = 1 \\
    x + 2y -z & = 3
  \end{aligned} \right. \Longrightarrow
  \left\{
  \begin{aligned}
    x + y + z & = 1 \\
    y - 2z & = -2
  \end{aligned} \right. \Longrightarrow 
  \left\{
    \begin{aligned}
    x + 3z & = -1 \\
    y - 2z & = 2
    \end{aligned} \right.
\end{align*}  

\hspace*{0.5cm} Solutionset: \[\{(x,y,z): x = -1 - 3t, y = 2 + 2t, z = t, t \in \mathbb{R}\}\]
\end{exercise}

\begin{remark}[title = Remark 1.3.2]
  The linear case is done by linear algebra: one can parameterize the solution set of a system of linear equations via Gaussian elimination. \\

  \hspace*{0.5cm}\textcolor{airforceblue}{\textit{(Go through the details yourself or otherwise ask.)}} 
\end{remark}

\begin{example}[title = Example 1.3.3 (Quadratic Case)]
  Let's consider everyone's favorite: the circle $x^2 + y^2 = 1$. 

  The circle has a rational parametrization: \[x = \frac{u(t)}{g(t)}, y = \frac{v(t)}{g(t)}, u, v, g \in \mathbb{R}[t]\]  
  Now consider the following:  

  \begin{minipage}{0.39\textwidth}
    \begin{center}
      \begin{tikzpicture}[scale=1.25]
        \def\t{0.65}
        
        % Achsen
        \draw[black!70, ->] (-1.5,0) -- (1.6,0);
        \draw[black!70, ->] (0,-1.2) -- (0,1.4);
        
        % Einheitskreis
        \draw[blue!70, thick] (0,0) circle (1);
        
        % Gerade durch (-1,0) mit y-Achsenabschnitt t
        \draw[kGreen!40!black, thick] (-1.3,{\t*(-1.3+1)}) -- (1.3,{\t*(1.3+1)});
        
        % Punkt (-1,0)
        \fill[kGreen!40!black] (-1,0) circle (1.2pt);
        
        % Klammer für t
        \draw[decorate, decoration={brace, amplitude=4pt, mirror}, green!40!black, thick]
        (0,0) -- (0,\t) node[midway, right=5pt] {$t$};
      \end{tikzpicture}
    \end{center}
  \end{minipage}
  \hspace*{0.2cm}
  \begin{minipage}{0.59\textwidth}
    The slope of the line can be changed, which is represented by the parameter $t$. Find a representation of the other intersection of the line with the circle (in rational parameterization).
  \end{minipage}
  \begin{center}
    \textcolor{gray}{------------------------------------------- 26.10.2026 --------------------------------------------} \\[0.3em]
  \end{center} 

  \paragraph{Solution via Breathing-Procedure:}\ \\
  First thing we want is an equation of the line: 
  % We know: y is related to x, t is the slope, we know p(0, -1) --> (x+1)  
  \[y = t \cdot (x+1)\]
  % y is becoming the slave of x because it is bind to it. 
  Combining both equations, we get: \[x^2 + (t \cdot (x+1))^2 = 1\]
  % Next Step: find $x$ in terms of $t$. 
  We know the term $(x-1)$ and have the following factorization: 
  \[(x + 1) \cdot (\ldots x + \ldots) = 0\] 
\end{example} 
\begin{example}[title = Example 1.3.3 (Quadratic Case) Continued]
  Now determine the other factor:
  \[(x + 1) \cdot ((1+t^2)\cdot x -1 + t^2) = 0\]
  For $x \neq 0$, we have \[x = \frac{1-t^2}{1+t^2} \text{ and } y = t(x+1) = t \left(\frac{1-t^2}{1+t^2} + 1\right) = \frac{2t}{1+t^2}\] 
  Double check: 
  \begin{align*}
    x^2 + y^2 & = \frac{(1-t^2)^2}{(1+t^2)^2} + \frac{4t^2}{(1+t^2)^2} \\
              & = \frac{1-2t^2+t^4 + 4t^2}{(1+t^2)^2}\\
              & = \frac{1+ 2t^2 + t^4}{(1 + t^2)^2} \\
              & = 1
  \end{align*}
\end{example} 

\begin{definition}[title = Definition 1.3.4 (Rational Functions)]
  A rational function with coefficients in the field $k$ in the variables $x_1, \ldots, x_n$ is a formal quotient $\frac{f}{g}$ where $f, g \in k[x_1, \dots, x_n]$ 
  and $g \neq 0$. The set of all rational functions is denoted by $k(x_1, \ldots, x_n)$ (round brackets). \\
  
  Two rational functions $\frac{f}{g}$ and $\frac{a}{b}$ are defined to be equal when $f \cdot b = a \cdot g$. \\

  Rational functions can be added and multiplied as follows:
\[\frac{f}{g} + \frac{u}{v} :=\frac{f \cdot v + g \cdot u}{g \cdot v}  \ \ \ \text{ and } \ \ \  \frac{f}{g} \cdot \frac{u}{v}:= \frac{f \cdot u}{g \cdot v} \]

  for $f, u, g, v \in k[x_1, \dots, x_n], g \neq 0, v \neq 0$\dots \\

  The above definitions are well-defined, because in the ring $k[x_1, \dots, x_n]$:
  \[g \neq 0 \text{ and } v \neq 0 \Rightarrow g \cdot v \neq 0\] 
  (This needs to be shown and it is not true in every ring.)
\end{definition}

\begin{remark}[title = Remark 1.3.5]
  $k[x_1, \dots, x_n]$ is a field (which is great!). 
\end{remark}

\begin{definition}[title = Definition 1.3.6]
  Let $V \subseteq k^n$ be a variety. We say that the expressions 
  \[
    x_1 = \frac{p_1(t_1, \ldots, t_m)}{q(t_1, \ldots, t_m)} , 
    \vdots , 
    x_n = \frac{p_n(t_1, \ldots, t_m)}{q(t_1, \ldots, t_m)}
  \]
  where $p_1, \ldots, p_n, q \in k[x_1, \dots, x_n] $ and $q \neq 0$ parameterize the variety $V$ if $V$ is the smallest variety w.r.t unclusion that contains the set 
  \[\left\{\left( \frac{p_1(t_1, \ldots, t_m)}{q(t_1, \ldots, t_m)}, \ldots, \frac{p_n(t_1, \ldots, t_m)}{q(t_1, \ldots, t_m)} \right): a_1, \ldots, a_m \in k, q(a_1, \ldots, a_m) \neq 0 \right\}\] 

  A variety that has such a parameterization (which is called rational) is called a \textit{rational variety}.
\end{definition} 

\begin{definition}[title = Definition 1.3.7]
  The description of a variety $V$ by a system of finitely many equations $f_1 = \ldots = f_m = 0$ where $f_1, \ldots, f_m \in k[x_1, \dots, x_n]$ is called an 
  \textit{implicit description} (this is our default description, it always exists by the definition of a variety).
  \begin{center}
    \scalebox{0.7}{%
      \begin{tikzpicture}[
        every node/.style={font=\itshape}, 
        blob/.style={draw=black!40!black, line width=1pt, ellipse,
                     minimum width=2.8cm, minimum height=1.5cm},
        arr/.style={draw=black!40!black, line width=1pt, -{Stealth[length=3mm]}}
      ]
        \node[blob] (imp) at (0,0) {implicit};
        \node[blob] (par) at (7,0) {parametric};
      
        \draw[arr] (imp.north east) to[bend left=25]
              node[midway, above] {parametrization} (par.north west);
        \draw[arr] (par.south west) to[bend left=25]
              node[midway, below] {implicitization} (imp.south east);
      
        % "not unique" note
        \draw[black!40!black, line width=1pt, double, double distance=1.5pt,
              -{Implies[length=2mm]}] (-1.6,-1.8) -- (-1.1,-0.9);
        \node[text=black!40!black] at (-1.6,-2.2) {not unique};
      \end{tikzpicture}}
  \end{center}
\end{definition} 

\begin{example}[title = Example 1.3.8]
  Bezier splines is a cursive given parameterically, whose $x$ and $y$ coordinates are parameterized as cubic polynomials. Normally written like this: 
  \[\begin{bmatrix}
    x \\ y
  \end{bmatrix} = (1 - t)^3 P_0 + 3(1 - t)^2 t P_1 + 3(1 - t) t^2 P_2 + t^3 P_3\] 
  where $P_0, \ldots, P_3 \in k^2$. 

  One may assk for an implicit equation of such a curve.  
    \begin{center}
      \scalebox{0.8}{%
        \begin{tikzpicture}[
          x={(0.8cm,-0.6cm)}, y={(1cm,0.5cm)}, z={(0cm,1cm)},
          every node/.style={font=\itshape},
          ax/.style={draw=green!40!black, line width=1pt, -{Stealth[length=3mm]}},
          curve/.style={draw=green!40!black, line width=1pt},
          proj/.style={draw=blue!60, line width=1pt},
          parr/.style={draw=blue!60, line width=1pt, -{Stealth[length=2.5mm]}}
        ]
          % Achsen
          \draw[ax] (0,0,0) -- (0,0,5.6) node[left] {t};
          \draw[ax] (0,0,0) -- (0,5,0) node[right] {y};
          \draw[ax] (0,0,0) -- (4.5,0,0) node[below right] {x};
        
          \draw[curve] plot[variable=\s, domain=-1.4:1.65, samples=80, smooth]
            ({2.5-\s*\s}, {1.8+\s*(\s*\s-1)}, {3-1.3*\s});
        
          \draw[proj] plot[variable=\s, domain=-1.4:1.4, samples=80, smooth]
            ({2.5-\s*\s}, {1.8+\s*(\s*\s-1)}, 0);
        
          \foreach \s in {-1,0,1,1.3} {
            \draw[parr] ({2.5-\s*\s}, {1.8+\s*(\s*\s-1)}, {3-1.3*\s})
                     -- ({2.5-\s*\s}, {1.8+\s*(\s*\s-1)}, 0);
          }
        \end{tikzpicture}
      }
    \end{center}
  Making things implicit $\leadsto$ projection
\end{example}

\begin{remark}[title = Remark 1.3.9]
  In $k^n$ one can define an affine subspace by a linear system as $\{x:Ax = b\}$ where $Ax = b$ is a linear system over $k$ and parametrically as $\{U \cdot t + v: t \in k^m\}$, where $U$ and $v$ are 
  a matrix reps. a vector over $k$. So, we have a linear version of our tasks. \\ 
  % TODo: bild 3  
  Both tasks are solvable with Gaussian Elimination.
\end{remark} 

\clearpage

\begin{exercise}[title = Exercise 1.3.10]
  \begin{enumerate}
    \item Why is the product of two non-zero polynomials in $k[x_1, \dots, x_n]$ a non-zero polynomial? 
    \paragraph{Solution:} 
    \item Find an implicit description of the parametric curve $x = \frac{t}{t+t^2}$ and $y = 1 - \frac{1}{t^2}$. 
    \paragraph{Solution:} 
    \item Justify Remark 1.3.9. 
    \paragraph{Solution:} 
    \item Find a rational representation of the sphere $V(x^2 + y^2 + z^2 - 1)$. 
    \paragraph{Solution:} 
    \item Find a parametric description of the variety $V(y^2 - x^2 + x^3)$. 
    \paragraph{Solution:} 
    \item Implicitize 
    \begin{align*}
      x_1 & = t_1 + t_2 \\
      x_2 & = t_1 + 2t_2 \\ 
      x_3 & = 2 t_1 + t_2 \\ 
      x_4 & = t_1 - t_2
    \end{align*} 
    Where $t_1, t_2$ are parameters.
    \paragraph{Solution:} 
   \end{enumerate}
   The field is $k = \mathbb{R}$. 
\end{exercise} 

